\documentclass[11pt,a4paper]{article}

% --- Packages ---
\usepackage[utf8]{inputenc}
\usepackage{amsmath, amssymb}
\usepackage{graphicx}
\usepackage{wrapfig}
\usepackage{multirow}
\usepackage{booktabs}
\usepackage{xcolor}
\usepackage{enumitem}
\usepackage{listings}
\usepackage[colorlinks=true, linkcolor=blue, citecolor=blue, urlcolor=blue]{hyperref}

% --- Listings configuration for Python ---
\lstset{
    language=Python,
    basicstyle=\ttfamily\small,
    keywordstyle=\color{blue}\bfseries,
    stringstyle=\color{red},
    commentstyle=\color{green!50!black},
    numbers=left,
    numberstyle=\tiny\color{gray},
    stepnumber=1,
    frame=single,
    breaklines=true,
    captionpos=b,
    showstringspaces=false
}

% --- Document Info ---
\title{Laplace Analysis and Applications}
\author{
    Pierre-Simon Laplace\textsuperscript{1}, 
    Oliver Heaviside\textsuperscript{2}, and 
    Gustav Robert Kirchhoff\textsuperscript{3}
}
\date{
    \small
    \textsuperscript{1}Académie des Sciences, Paris, France \\
    \textsuperscript{2}Royal Society, London, United Kingdom \\
    \textsuperscript{3}University of Heidelberg, Heidelberg, Germany
}

\begin{document}

\maketitle

\section{From Functions to Transforms}

\subsection{The Laplace Transform}

Let $f(t)$ be a function defined for $t \ge 0$. Its Laplace transform is defined by
\begin{equation}
    \mathcal{L}\{f(t)\} = F(s) = \int_{0}^{\infty} e^{-st}f(t)dt, \quad s>0 \label{eq:laplace_def}
\end{equation}

The transform changes a function of time into a function of the new variable $s$.\footnote{A concise overview is available from the Laplace Transform entry on \href{https://mathworld.wolfram.com/LaplaceTransform.html}{MathWorld}.} Symbolically, $f(t) \rightarrow F(s)$. The inverse Laplace transform reverses this process:
\begin{equation}
    f(t) = \mathcal{L}^{-1}\{F(s)\}.
\end{equation}

The main purpose of the method is to replace differentiation and integration in the time domain by simpler algebraic operations in the transform domain \cite{kreyszig}. 

\subsection{A Simple Example}

For the constant function $f(t)=1$, 
\begin{equation}
    \mathcal{L}\{1\} = \int_{0}^{\infty} e^{-st}dt = \left[ -\frac{1}{s}e^{-st} \right]_{0}^{\infty} = \frac{1}{s}
\end{equation}
Similarly,
\begin{equation}
    \mathcal{L}\{t\} = \frac{1}{s^{2}}, \quad \mathcal{L}\{t^{2}\} = \frac{2}{s^{3}}
\end{equation}
The variable $s$ is commonly interpreted as a complex-frequency variable.

\section{Common Transform Equations}

Table \ref{tab:transforms} lists several frequently used Laplace-transform equations.

\begin{table}[htpb]
    \centering
    \caption{Selected Laplace-transform equations.}
    \label{tab:transforms}
    \renewcommand{\arraystretch}{1.5}
    \begin{tabular}{|c|c|c|c|}
        \toprule
        \multirow{2}{*}{\textbf{Family}} & \multirow{2}{*}{\textbf{Time function}} & \multicolumn{2}{c}{\textbf{Laplace transform}} \\
        \cmidrule{3-4}
        & & Equation & Condition \\
        \midrule
        \multirow{3}{*}{Algebraic} & $1$ & $1/s$ & $s>0$ \\
        & $t$ & $1/s^2$ & $s>0$ \\
        & $t^n$ & $\frac{n!}{s^{n+1}}$ & $s>0, n=0,1,2,\dots$ \\
        \midrule
        \multirow{2}{*}{Exponential} & $e^{at}$ & $\frac{1}{s-a}$ & $s>a$ \\
        & $e^{-at}$ & $\frac{1}{s+a}$ & $s>-a$ \\
        \midrule
        \multirow{2}{*}{Oscillatory} & $\sin(\omega t)$ & $\frac{\omega}{s^2+\omega^2}$ & $\omega>0$ \\
        & $\cos(\omega t)$ & $\frac{s}{s^2+\omega^2}$ & $\omega>0$ \\
        \midrule
        & Unit-step function $u(t-a)$ & $\frac{e^{-as}}{s}$ & $a>0$ \\
        \bottomrule
    \end{tabular}
\end{table}

The formulas in Table \ref{tab:transforms} allow many functions to be transformed without repeatedly evaluating improper integrals.

\section{Application to an RC Circuit}

\begin{wrapfigure}{r}{0.4\textwidth}
    \centering
    % The instructions requested a rotated figure. The 'angle=90' parameter fulfills this requirement.
    % Make sure a file named 'circuit.png' exists in your working directory.
    \includegraphics[width=\linewidth]{rc_circuit.png}
    \caption{A simple series RC circuit.}
    \label{fig:circuit}
\end{wrapfigure}

\textbf{Circuit Model} \\
Figure \ref{fig:circuit} shows a simple resistor--capacitor circuit connected to an input voltage source. The resistor limits the current, while the capacitor stores electrical energy. 

For a resistor $R$ and capacitor $C$, Kirchhoff's voltage law gives
\begin{equation}
    RC\frac{dv_{C}(t)}{dt} + v_{C}(t) = v_{in}(t), \label{eq:kvl}
\end{equation}
where $v_{C}(t)$ is the voltage across the capacitor. Assuming $v_{C}(0)=0$, taking the Laplace transform of Equation (\ref{eq:kvl}) gives
\begin{align}
    RC[sV_{C}(s) - v_{C}(0)] + V_{C}(s) &= V_{in}(s) \nonumber \\
    \Rightarrow RCsV_{C}(s) + V_{C}(s) &= V_{in}(s) \nonumber \\
    \Rightarrow V_{C}(s)(RCs + 1) &= V_{in}(s). \label{eq:laplace_rc}
\end{align}

The transfer function is therefore
\begin{equation}
    H(s) = \frac{V_{C}(s)}{V_{in}(s)} = \frac{1}{RCs+1}.
\end{equation}
For a constant input $V_0$, $V_{in}(s) = \frac{V_0}{s}$. Thus,
\begin{equation}
    V_{C}(s) = \frac{V_0}{s(RCs+1)}.
\end{equation}
Taking the inverse transform gives
\begin{equation}
    v_{C}(t) = V_0(1 - e^{-t/(RC)}). \label{eq:time_domain}
\end{equation}

The quantity $\tau = RC$ is called the time constant of the circuit. A larger value of $\tau$ causes the capacitor to charge more slowly \cite{ogata}.

\section{How Laplace Analysis Is Used}

The Laplace transform is useful in many areas:
\begin{description}
    \item[\textcolor{blue}{Differential equations}] initial conditions are included automatically; derivatives become algebraic expressions.
    \item[\textcolor{blue}{Electrical circuits}] voltages and currents are studied in the $s$-domain; components lead to transfer functions.
    \item[\textcolor{blue}{Control and signal systems}] stability can be studied through poles; common test inputs include:
    \begin{itemize}
        \item the unit-step input;
        \item the impulse input;
        \item sinusoidal inputs.
    \end{itemize}
\end{description}

In the time domain, a system may be described by a differential equation, whereas in the Laplace domain, the same system is often described by a rational algebraic expression.

\section{A Short Python Implementation}

Listing \ref{lst:rc_code} evaluates the capacitor voltage at a finite set of time instants using an external Python file.

\begin{lstlisting}[caption={Computing samples of the RC step response.}, label={lst:rc_code}]
from math import exp

def rc_response(resistance, capacitance, input_voltage, times):
    """Return capacitor-voltage samples for an RC step input."""
    if resistance <= 0 or capacitance <= 0:
        raise ValueError("R and C must be positive")
        
    tau = resistance * capacitance
    values = []
    
    for time in times:
        voltage = input_voltage * (1 - exp(-time / tau))
        values.append(voltage)
        
    return values

times = [0.0, 0.5, 1.0, 2.0]
print(rc_response(2.0, 0.5, 5.0, times))
\end{lstlisting}

Equation (\ref{eq:time_domain}) and Listing \ref{lst:rc_code} describe the same computation.

\begin{thebibliography}{9}
\bibitem{kreyszig} 
E. Kreyszig, \textit{Advanced Engineering Mathematics}, 10th ed. John Wiley \& Sons, 2011.

\bibitem{ogata} 
K. Ogata, \textit{Modern Control Engineering}, 5th ed. Prentice Hall, 2010.
\end{thebibliography}

\end{document}